\documentclass[10pt,a4paper]{amsart}
\usepackage{amsmath,amssymb,amsthm,graphicx}
% Retained minimal preamble for the Stacks Project chapter more-algebra.tex (More on Algebra).
% The chapter file pinned-original-source.tex carries only \input{preamble} on its first line, so no
% preamble text lives inside the chapter itself.  The Stacks Project uses one single shared preamble for
% every chapter; the snapshot stages it identically as sources/stacks-chow/preamble.tex and
% sources/stacks-etale-cohomology/preamble.tex.  The block below is copied mechanically, line for line,
% from sources/stacks-chow/preamble.tex lines 14-18 (the Xy-pic setup) and 154-217 (the native theorem
% declarations together with the mathematical macros).  Nothing was added, removed or redefined; the
% wrapper class and the amsmath/amssymb/amsthm/graphicx base packages are supplied by the assembler.
% Chapter: more-algebra.tex -- More on Algebra.
% For commutative diagrams we use Xy-pic
\usepackage[all]{xy}

% We use 2cell for 2-commutative diagrams.
\xyoption{2cell}

% Theorem environments.
%
\theoremstyle{plain}
\newtheorem{theorem}[subsection]{Theorem}
\newtheorem{proposition}[subsection]{Proposition}
\newtheorem{lemma}[subsection]{Lemma}

\theoremstyle{definition}
\newtheorem{definition}[subsection]{Definition}
\newtheorem{example}[subsection]{Example}
\newtheorem{exercise}[subsection]{Exercise}
\newtheorem{situation}[subsection]{Situation}

\theoremstyle{remark}
\newtheorem{remark}[subsection]{Remark}
\newtheorem{remarks}[subsection]{Remarks}

\numberwithin{equation}{subsection}

% Macros
%
\def\lim{\mathop{\mathrm{lim}}\nolimits}
\def\colim{\mathop{\mathrm{colim}}\nolimits}
\def\Spec{\mathop{\mathrm{Spec}}}
\def\Hom{\mathop{\mathrm{Hom}}\nolimits}
\def\Ext{\mathop{\mathrm{Ext}}\nolimits}
\def\SheafHom{\mathop{\mathcal{H}\!\mathit{om}}\nolimits}
\def\SheafExt{\mathop{\mathcal{E}\!\mathit{xt}}\nolimits}
\def\Sch{\mathit{Sch}}
\def\Mor{\mathop{\mathrm{Mor}}\nolimits}
\def\Ob{\mathop{\mathrm{Ob}}\nolimits}
\def\Sh{\mathop{\mathit{Sh}}\nolimits}
\def\NL{\mathop{N\!L}\nolimits}
\def\CH{\mathop{\mathrm{CH}}\nolimits}
\def\proetale{{pro\text{-}\acute{e}tale}}
\def\etale{{\acute{e}tale}}
\def\QCoh{\mathit{QCoh}}
\def\Ker{\mathop{\mathrm{Ker}}}
\def\Im{\mathop{\mathrm{Im}}}
\def\Coker{\mathop{\mathrm{Coker}}}
\def\Coim{\mathop{\mathrm{Coim}}}

% Boxtimes
%
\DeclareMathSymbol{\boxtimes}{\mathbin}{AMSa}{"02}

%
% Macros for moduli stacks/spaces
%
\def\QCohstack{\mathcal{QC}\!\mathit{oh}}
\def\Cohstack{\mathcal{C}\!\mathit{oh}}
\def\Spacesstack{\mathcal{S}\!\mathit{paces}}
\def\Quotfunctor{\mathrm{Quot}}
\def\Hilbfunctor{\mathrm{Hilb}}
\def\Curvesstack{\mathcal{C}\!\mathit{urves}}
\def\Polarizedstack{\mathcal{P}\!\mathit{olarized}}
\def\Complexesstack{\mathcal{C}\!\mathit{omplexes}}
% \Pic is the operator that assigns to X its picard group, usage \Pic(X)
% \Picardstack_{X/B} denotes the Picard stack of X over B
% \Picardfunctor_{X/B} denotes the Picard functor of X over B
\def\Pic{\mathop{\mathrm{Pic}}\nolimits}
\def\Picardstack{\mathcal{P}\!\mathit{ic}}
\def\Picardfunctor{\mathrm{Pic}}
\def\Deformationcategory{\mathcal{D}\!\mathit{ef}}

% --- appended 2026-10-12 (Chief): the project own definitions of the reference/slogan/history
% display environments, copied mechanically from the shared Stacks preamble lines 8-12
% (\usepackage{verbatim} plus the three \newenvironment lines, which the chapter preamble keeps
% verbatim from the project preamble). Nothing added, removed or redefined: these environments
% carry no mathematical content and the project itself typesets them as comments. Without these
% three lines a counted statement whose own text uses them fails with Environment slogan undefined.
\usepackage{verbatim}
\newenvironment{reference}{\comment}{\endcomment}
\newenvironment{slogan}{\comment}{\endcomment}
\newenvironment{history}{\comment}{\endcomment}
\usepackage{cleveref}
\crefname{originallemma}{original lemma}{original lemmas}
\Crefname{originallemma}{Original lemma}{Original lemmas}
\makeatletter
\newlabel{lemma-finite-push-pseudo-coherent}{{064Z}{source}{}{sourcepointer-1}{}}
\newlabel{lemma-finite-push-pseudo-coherent@cref}{{[originallemma][0][]064Z}{[1][0][]source}{}{}{}}
\makeatother
\crefname{originallemma}{original lemma}{original lemmas}
\Crefname{originallemma}{Original lemma}{Original lemmas}
\makeatletter
\newlabel{lemma-pull-pseudo-coherent}{{0650}{source}{}{sourcepointer-2}{}}
\newlabel{lemma-pull-pseudo-coherent@cref}{{[originallemma][0][]0650}{[1][0][]source}{}{}{}}
\makeatother
\title{Relative pseudo-coherence survives localizing the algebra and the base ring at once}
\author{The Stacks Project Authors}
\date{stacks-project:more-algebra.tex:local-snapshot-20261012; GNU FDL 1.2 or later}
\begin{document}
\maketitle
\setcounter{subsection}{1}
\noindent\textit{Editorial note.} All delivered mathematical text is the author's own, copied from the pinned Stacks Project chapter file more-algebra.tex; no statement or proof step was rewritten, reordered or removed. The counted statement is the labeled lemma environment at source lines 21498-21505 (canonical tag 0679) and the counted proof is the proof environment at source lines 21507-21540, which immediately follows it and is its own proof.\par\medskip
\noindent\textit{Editorial note.} The printed text begins a sentence with a lower case write at proof line 21531 and writes A\_f for a localization where A\_g is meant at proof lines 21527-21529; both are copied exactly as printed and are not corrected.\par\medskip
\noindent\textit{Editorial note.} The preamble span runs to line 90, the current end of the retained preamble, which includes the appended environment definitions; the counted text does not use them, and the span is taken whole so that the file version is unambiguous.\par\medskip
\noindent\textit{Editorial note.} The source URL is the source-level Stacks Project URL, as required when more than one card is produced from one source; the per-claim tag is carried by claim\_key and by this note. The counted statement and proof cite no bibliography entry, so this card needs no bibliography.\par\medskip
\noindent\textit{Allowed standard prerequisites (editorial).}\begin{itemize}
\item A complex that is m-pseudo-coherent relative to a base remains so when it is viewed as a complex over a larger polynomial algebra over the same base, provided the algebra used is pseudo-coherent as a module over that larger algebra. invoked at proof lines 21515-21518 and again at proof lines 21535-21539. Moves relative pseudo-coherence from the localized base to a polynomial ring over the base and then to the quotient presenting the further localization.
\item The localization of a ring at an element is the quotient of a polynomial ring by the relation expressing the inverse of that element, and the same holds for a second localization. used at proof lines 21509-21511 and 21531-21534. Provides the explicit presentations in terms of which the two moves are made.
\item Pulling back a relatively pseudo-coherent complex along a ring map of finite presentation preserves relative pseudo-coherence. invoked at proof lines 21521-21530. Identifies the tensor product of the complex with the localization of the polynomial algebra with the complex over the localized algebra.
\end{itemize}\medskip
\ifdefined\hypertarget\hypertarget{sourcepointer-1}{}\fi
\noindent\textit{Source reference.} Original lemma 064Z: sources/stacks-more-algebra/pinned-original-source.tex, source lines 16283--16283. The reference is applied twice, at proof lines 21516-21518 and at proof lines 21536-21538, to move relative pseudo-coherence to the polynomial algebra over the base; the presentations that allow it are inside the counted span.\par\medskip
\ifdefined\hypertarget\hypertarget{sourcepointer-2}{}\fi
\noindent\textit{Source reference.} Original lemma 0650: sources/stacks-more-algebra/pinned-original-source.tex, source lines 16345--16345. The reference is applied at proof lines 21521-21527 to identify the tensor product of the complex with the localization of the polynomial algebra with the complex over the localized algebra; the second transfer built on that is inside the counted span.\par\medskip
\begin{lemma}
\label{lemma-localize-relative-pseudo-coherent}
Let $R$ be a ring, $f \in R$ an element, $R_f \to A$ is a finite type ring map,
$g \in A$, and $K^\bullet$ a complex of $A$-modules.
If $K^\bullet$ is $m$-pseudo-coherent (resp.\ pseudo-coherent)
relative to $R_f$, then $K^\bullet \otimes_A A_g$ is
$m$-pseudo-coherent (resp.\ pseudo-coherent) relative to $R$.
\end{lemma}

\begin{proof}
First we show that $K^\bullet$ is $m$-pseudo-coherent relative to $R$.
Namely, suppose $R_f[x_1, \ldots, x_n] \to A$ is surjective. Write
$R_f = R[x_0]/(fx_0 - 1)$. Then $R[x_0, x_1, \ldots, x_n] \to A$
is surjective, and $R_f[x_1, \ldots, x_n]$ is pseudo-coherent as
an $R[x_0, \ldots, x_n]$-module. Hence by
Lemma \ref{lemma-finite-push-pseudo-coherent}
we see that $K^\bullet$ is $m$-pseudo-coherent as a complex of
$R[x_0, x_1, \ldots, x_n]$-modules.

\medskip\noindent
Choose an element $g' \in R[x_0, x_1, \ldots, x_n]$ which maps to $g \in A$. By
Lemma \ref{lemma-pull-pseudo-coherent}
we see that
\begin{align*}
K^\bullet \otimes_{R[x_0, x_1, \ldots, x_n]}^{\mathbf{L}}
R[x_0, x_1, \ldots, x_n, \frac{1}{g'}] & =
K^\bullet \otimes_{R[x_0, x_1, \ldots, x_n]}
R[x_0, x_1, \ldots, x_n, \frac{1}{g'}] \\
& = K^\bullet \otimes_A A_f
\end{align*}
is $m$-pseudo-coherent as a complex of
$R[x_0, x_1, \ldots, x_n, \frac{1}{g'}]$-modules.
write
$$
R[x_0, x_1, \ldots, x_n, \frac{1}{g'}] =
R[x_0, \ldots, x_n, x_{n + 1}]/(x_{n + 1}g' - 1).
$$
As $R[x_0, x_1, \ldots, x_n, \frac{1}{g'}]$ is pseudo-coherent as a
$R[x_0, \ldots, x_n, x_{n + 1}]$-module we conclude (see
Lemma \ref{lemma-finite-push-pseudo-coherent})
that $K^\bullet \otimes_A A_g$ is $m$-pseudo-coherent as a complex of
$R[x_0, \ldots, x_n, x_{n + 1}]$-modules as desired.
\end{proof}

\end{document}
